\documentclass[../00_Main.tex]{subfiles}

\begin{document}

Electroencephalography (\gls{eeg}) is a widely used, non-invasive technique for measuring electrical activity generated by neural populations in the brain.
Its applications span clinical diagnostics, neuroscience research, and the development of brain--computer interfaces (\gls{bci}). As \gls{eeg} systems evolve toward
higher channel counts, wearable form factors, and tighter integration with digital signal processing pipelines, the need for reliable and reproducible validation
methodologies becomes increasingly critical.

A fundamental challenge in \gls{eeg} system evaluation lies in the absence of a known ground truth when measurements are performed on human subjects. Biological
variability, uncontrollable physiological states, and environmental noise make it difficult to isolate the performance of individual system components---such
as electrodes, acquisition hardware, or signal-processing algorithms---from the underlying neural activity. As a result, controlled testing and quantitative
benchmarking of \gls{eeg} systems remain inherently limited when relying exclusively on in-vivo experiments.

Physical head phantoms provide a controlled alternative for \gls{eeg} validation. By replicating the electrical behavior of the head and enabling the injection
of predefined signals, phantoms allow repeatable experiments under well-defined conditions. However, existing \gls{eeg} phantom solutions exhibit several practical
limitations. Many rely on specialized or proprietary materials or involve complex or expensive manufacturing processes, which can hinder accessibility and adaptation across research contexts.

A further limitation of current approaches is the frequent separation between physical experimentation, numerical modeling, and signal-generation software. Phantom heads are often treated as isolated test objects, without an explicit linkage between injected signals, material properties, and measured surface potentials. This separation complicates quantitative interpretation and limits the systematic tuning of phantom-based experiments.

Within this scope, the work develops the following components:

\begin{enumerate}
    \item \textbf{A parametric \gls{cad} model of a phantom head}, defining geometry, electrode placement, and internal routing suitable for controlled signal injection and measurement. The model provides a complete, printable representation of the phantom and establishes a consistent geometric reference for subsequent experimental and numerical work.
    \item \textbf{An experimentally characterized conductive filling medium based on a saline--gelatine mixture}, intended as the filling material of the phantom head. The electrical properties of the medium are evaluated across the relevant frequency range to verify effective signal transmission, demonstrate tunability of conductivity through ionic concentration of the solute, and assess reproducibility across preparations, supporting its use for controlled transmission of \gls{eeg}-like signals.
    \item \textbf{A signal generation and acquisition platform}, combining a cross-platform desktop application with dedicated signal-generation hardware. The platform enables deterministic multi-channel excitation, real-time visualization, structured data logging, and integration with OpenBCI acquisition systems for closed-loop validation.
    %\\ The software platform is designed around a set of functional and usability requirements: comprehensive waveform synthesis capabilities spanning Delta, Theta, Alpha, and Beta frequency bands characteristic of brain activity; parametric control allowing real-time modification of frequency and amplitude with low latency and jitter (within acceptable bounds of 100ms and 20ms respectively); and a user-friendly interface with intuitive controls that enable researchers to configure experiments without extensive training or technical knowledge of the underlying implementation.
    \item \textbf{A conditioned forward and inverse system response of the phantom head}, providing a calibrated, control-oriented model linking internal actuation points to surface electrode measurements. This system response is derived through numerical modeling of volumetric electrical conduction within the phantom head, incorporating: the geometry defined in (1), the internal conductive medium characterized in (2), mocking the excitation signals that would be generated by the proposed software application (3). \\The raw system response undergoes numerical conditioning to ensure stability and reconstruction suitability, yielding a bidirectional design tool that supports both forward mappings (predicting surface measurements from actuation patterns) and inverse mappings (determining internal excitation configurations from target surface potentials). This conditioned response enables users to compute the required input signals for the application (3) to produce desired surface electrode measurements, or conversely, to estimate expected surface measurements at any electrode location given any input signal configuration—facilitating controlled, reproducible experimental configurations with predictable outcomes.
\end{enumerate}

These four components are designed to operate as a single integrated tool, where the physical phantom (1, 2) provides a controlled test environment, the software platform (3) enables systematic signal injection and measurement visualization, and the numerical model (4) serves as a predictive bridge between experimental design and measured outcomes. This integration is designed with modularity in mind: the software platform (3) remains independent of any specific phantom geometry, allowing users to adapt component (4) to different physical configurations while maintaining a consistent signal generation and visualization interface, providing flexibility for various use cases and prototype testing.


\nota{PODRIAMOS PONER ESTO EN EL PUNTO 3? CREO QUE QUEDARIA MEJOR ASI NO SUENA A QUE VOLVIMOS AL MISMO TEMA, SINO QUE EXPLICAMOS TODOS LOS COMPONENTES Y DESPUES CERRAMOS CON LA INTEGRACION DE LOS 4 COMO UNA SOLA HERRAMIENTA, QUE A SU VEZ FUE PENSADA DE FORMA QUE PUEDE SER MODULAR, Y POR LO TANTO ADAPTABLE A DIFERENTES CONFIGURACIONES FISICA.  DEJE UNA VERSION  COMENTADA EN EL ITEM 3 DE COMO PODRIAMOS PONERLO, SI PODES CHEQUEALA }

The software platform design is based on a set of requirements. First, the system must support comprehensive waveform synthesis capabilities, including the generation of Delta, Theta, Alpha, and Beta frequency bands that are characteristic of brain activity. To enable fine-tuned experimentation, the system should provide parametric control that allows real-time modification of both frequency (measured in Hz) and amplitude of the generated signals, maintaining low latency and jitter within acceptable bounds of $100\text{ms}$ and $20\text{ms}$ respectively. Last but not least, usability considerations demand a user-friendly interface design with intuitive controls that enable researchers to quickly configure experiments without extensive training or technical knowledge of the underlying implementation.

\biblio % Needed for referencing when compiling individual subfiles - Do not remove
\end{document}